\chapter{Leg Pain}

\section*{Definition}
Pain or discomfort anywhere in the lower extremity from the hip to the foot, arising from musculoskeletal, vascular, neurological, or dermatological sources.

\section*{Pathophysiology}
Leg pain mechanisms include nociceptive (musculoskeletal injury, inflammation), neuropathic (nerve compression or damage from disc herniation, peripheral neuropathy), ischemic (peripheral arterial disease, DVT), and referred (hip or spine pathology presenting as thigh or leg pain).

\section*{Etiology}
\begin{itemize}
  \item Peripheral arterial disease (claudication)
  \item Deep vein thrombosis (DVT)
  \item Sciatica or lumbar radiculopathy
  \item Muscle strain or tear (hamstring, quadriceps, calf)
  \item Peripheral neuropathy (diabetic, alcoholic)
  \item Varicose veins or chronic venous insufficiency
  \item Compartment syndrome (acute or chronic)
  \item Shin splints (medial tibial stress syndrome)
  \item Stress fracture
  \item Cellulitis or soft tissue infection
\end{itemize}

\section*{Indications for Evaluation}
Seek care if:
\begin{itemize}
  \item Severe or worsening symptoms
  \item Sudden severe leg pain with swelling and redness (possible DVT), leg pain with chest pain or shortness of breath, or non-healing ulcers
  \item Accompanied by other concerning signs
\end{itemize}

\section*{Related Symptoms}
\begin{itemize}
  \item Swelling
  \item Numbness or tingling
  \item Skin changes
  \item Ulcers
  \item Calf tenderness
\end{itemize}
\textbf{Note:} Always consider serious underlying conditions when evaluating persistent leg pain.

\section*{Self-Care / Home Management}
\begin{itemize}
  \item Rest and avoid activities that may aggravate the condition.
  \item Maintain adequate hydration and a balanced diet.
  \item Over-the-counter remedies may provide symptomatic relief (consult a pharmacist or doctor).
  \item Monitor symptoms and seek professional advice if they persist or worsen.
  \item Avoid self-medicating with prescription medications without medical guidance.
\end{itemize}

\section*{Diagnostic Approach}
\begin{itemize}
  \item A thorough history and physical examination are the first steps in evaluation.
  \item Laboratory tests (blood work, urinalysis) may be ordered based on clinical suspicion.
  \item Imaging studies (X-ray, ultrasound, CT, MRI) may be indicated when structural pathology is suspected.
  \item Specialist referral is arranged when the diagnosis remains uncertain or requires specific expertise.
\end{itemize}

\section*{Treatment / Management Overview}
\begin{itemize}
  \item Treatment targets the underlying cause identified through diagnostic evaluation.
  \item Supportive care and symptom management are provided while awaiting definitive therapy.
  \item Medications, lifestyle modifications, and physical therapies may be prescribed as appropriate.
  \item Follow-up ensures treatment effectiveness and allows timely adjustments to the care plan.
\end{itemize}

\clearpage